% Volume VII: Deep Learning Mathematics
% Continuity note for Chapter 38.
% Previous dependencies: Chapter 7 establishes shape-first gradients, Jacobians, Frobenius pairings, and matrix calculus; Chapters 20-23 establish optimization, computational graphs, forward/reverse-mode autodiff, vector-Jacobian products, and numerical stability; Chapter 37 defines feedforward networks as typed compositions.
% New durable concepts: chain rule for networks; computational graphs for layers; reverse-mode accumulation; adjoints; vector-Jacobian products; layerwise backpropagation; parameter gradients for weights and biases; batch gradients; differentiable programming; differentiable simulators and control flow; custom gradients; non-differentiable operations; vanishing and exploding gradients; saturation; dead ReLUs; ill-conditioning; residual paths; clipping; gradient debugging.
% Forward hooks: Chapter 39 studies deep optimization and training dynamics; Chapter 40 applies backpropagation to CNNs, RNNs, transformers, and GNNs; Chapter 41 differentiates generative-model objectives; Chapters 44-51 reuse differentiable simulations and likelihood gradients for neuroscience models.
% Notation commitments: Scalar losses are L; Jacobians use numerator layout J_{ij}=partial y_i/partial x_j; reverse adjoints are barred variables such as \bar h_\ell=\nabla_{h_\ell}L; layer preactivations are z_\ell; activations are h_\ell; elementwise products use \odot; scalar-loss matrix gradients satisfy dL=<nabla_W L,dW>_F.
\section{Chapter 38: Backpropagation and Differentiable Programming}
\label{ch:backpropagation-differentiable-programming}

Chapter 37 defined neural networks as compositions of typed maps. To train such a network, we need derivatives of a scalar loss with respect to many parameters. Chapter 7 taught that these derivatives have shapes and are defined through differentials and Frobenius pairings. Chapter 23 taught that reverse-mode automatic differentiation computes those derivatives on computational graphs. This chapter specializes those ideas to neural networks and then broadens them again to differentiable programs.

Backpropagation is not a biologically complete theory of learning. It is an efficient reverse-mode chain-rule algorithm for parameterized computations. Its power comes from reusing local derivatives and cached forward values. Its failure modes come from the same chain rule: long products of Jacobians can shrink, explode, saturate, or become numerically unstable.

\paragraph{Chapter problem.}
How does reverse-mode differentiation compute shape-correct gradients through neural networks and general differentiable programs, and why can gradient flow fail?

\paragraph{Section map.}
Section 38.1 develops the chain rule for networks using Jacobians and vector-Jacobian products. Section 38.2 derives the backpropagation algorithm and weight/bias gradients. Section 38.3 generalizes backpropagation to differentiable programming with arbitrary differentiable modules and control flow. Section 38.4 analyzes vanishing, exploding, saturated, dead, and ill-conditioned gradients.

\subsection{38.1 Chain rule for networks}

\paragraph{Guiding question.}
How do derivatives move backward through a composed network without forming every full Jacobian explicitly?

\paragraph{Beginner-facing intuition.}
If \(L\) depends on \(u\), and \(u\) depends on \(x\), then changing \(x\) changes \(u\), which changes \(L\). That is the chain rule. In a neural network, this happens through many vector-valued intermediate variables. The local derivative of one layer is a Jacobian, but backpropagation usually multiplies by Jacobians rather than storing them.

The key object is a vector-Jacobian product. If \(y=f(x)\) and \(L=L(y)\), then the upstream gradient \(\bar y=\nabla_y L\) is pulled back to
\[
\bar x=J_f(x)^\top \bar y.
\]
This has the same shape as \(x\).

\paragraph{Formal definitions.}
Let
\[
f:\mathbb R^n\to\mathbb R^m
\]
be differentiable at \(x\), with Jacobian
\[
J_f(x)\in\mathbb R^{m\times n},
\qquad
(J_f)_{ij}=\frac{\partial f_i}{\partial x_j}.
\]
Let \(L:\mathbb R^m\to\mathbb R\) be a scalar loss. Define
\[
y=f(x),\qquad \bar y=\nabla_y L(y)\in\mathbb R^m.
\]
The chain rule gives
\[
\boxed{
\bar x=\nabla_x L(f(x))=J_f(x)^\top \bar y\in\mathbb R^n.
}
\]
This is a vector-Jacobian product written in column-gradient convention.

For a composition
\[
x_0\xrightarrow{f_1}x_1\xrightarrow{f_2}\cdots\xrightarrow{f_L}x_L\xrightarrow{}L,
\]
the reverse recursion is
\[
\boxed{
\bar x_{\ell-1}=J_{f_\ell}(x_{\ell-1})^\top \bar x_\ell.
}
\]

\paragraph{Core derivation from differentials.}
Let \(y=f(x)\). A small perturbation \(dx\) produces
\[
dy=J_f(x)\,dx.
\]
Because \(L\) is scalar,
\[
dL=\bar y^\top dy
=\bar y^\top J_f(x)\,dx
=(J_f(x)^\top\bar y)^\top dx.
\]
By the definition of the gradient as the vector satisfying \(dL=\bar x^\top dx\), we identify
\[
\bar x=J_f(x)^\top\bar y.
\]
This is exactly the shape-first matrix-calculus rule from Chapter 7.

\paragraph{Concrete example.}
Let
\[
f(x)=
\begin{pmatrix}
x_1+x_2\\
x_1x_2
\end{pmatrix},
\qquad
L(y)=y_1+2y_2.
\]
Then
\[
J_f(x)=
\begin{pmatrix}
1&1\\
x_2&x_1
\end{pmatrix},
\qquad
\bar y=\nabla_y L=
\begin{pmatrix}1\\2\end{pmatrix}.
\]
At \(x=(3,4)^\top\),
\[
\bar x=J_f^\top\bar y
=
\begin{pmatrix}
1&4\\
1&3
\end{pmatrix}
\begin{pmatrix}1\\2\end{pmatrix}
=
\begin{pmatrix}9\\7\end{pmatrix}.
\]
Directly, \(L=x_1+x_2+2x_1x_2\), so \(\nabla_xL=(1+2x_2,1+2x_1)^\top=(9,7)^\top\).

\paragraph{Computational neuroscience example.}
In a differentiable encoding model, stimulus features \(x\) pass through filters and nonlinearities to predict firing rate \(r\), and a loss compares \(r\) to observed spikes. The vector-Jacobian product tells how a change in stimulus feature, filter weight, or latent state would change the loss. The derivative is a property of the model, not direct evidence for how the biological circuit changes synapses.

\paragraph{AI and machine-learning example.}
In a classifier, the softmax-cross-entropy loss produces an upstream gradient at the logits. Backpropagation pulls that gradient through linear layers, activation functions, attention blocks, and embeddings by repeated vector-Jacobian products. This is why a scalar training loss can update millions or billions of parameters.

\paragraph{Common misconceptions and failure modes.}
The chain rule for vectors is not ordinary scalar multiplication; the order and transpose of Jacobians matter. Gradients have the same shape as the variables they differentiate. Reverse mode does not require forming a full Jacobian when a vector-Jacobian product can be computed directly. Reusing a variable in two downstream branches means gradient contributions add.

\paragraph{Exercises.}
\begin{enumerate}[label=\textbf{38.1.\arabic*.}, leftmargin=*]
\item \textbf{Mechanical.} If \(f:\mathbb R^5\to\mathbb R^3\), state the shape of \(J_f\), \(\bar y\), and \(J_f^\top\bar y\).
\item \textbf{Proof or derivation.} Use differentials to derive \(\nabla_x L(f(x))=J_f(x)^\top\nabla_y L\).
\item \textbf{Numerical/computational.} For the example \(f(x)=(x_1+x_2,x_1x_2)\), compute \(\bar x\) at \(x=(1,5)\) with the same loss.
\item \textbf{Interpretation.} What does \(\bar h_\ell=\nabla_{h_\ell}L\) mean for a hidden representation in a neural network?
\item \textbf{Misconception/debugging.} A student writes \(\bar x=J_f\bar y\) instead of \(J_f^\top\bar y\). Use shapes to diagnose the mistake.
\end{enumerate}

\subsection{38.2 Backpropagation algorithm}

\paragraph{Guiding question.}
For a feedforward network, what exactly is computed in the forward pass and backward pass?

\paragraph{Beginner-facing intuition.}
Backpropagation has two phases. The forward pass computes and caches intermediate values: preactivations, activations, and the loss. The backward pass starts from the loss derivative at the output and moves backward layer by layer. At each layer it computes three things: the gradient with respect to the weights, the gradient with respect to the bias, and the gradient to pass to the previous layer.

The formulas are short because an affine layer has simple differentials.

\paragraph{Formal definitions and notation.}
For one training example, let
\[
h_0=x,
\qquad
z_\ell=W_\ell h_{\ell-1}+b_\ell,
\qquad
h_\ell=\phi_\ell(z_\ell),
\]
where
\[
W_\ell\in\mathbb R^{d_\ell\times d_{\ell-1}},
\quad
b_\ell\in\mathbb R^{d_\ell},
\quad
h_\ell,z_\ell\in\mathbb R^{d_\ell}.
\]
Let \(L=L(h_L,y)\) be a scalar loss. Define adjoints
\[
\bar h_\ell=\nabla_{h_\ell}L,\qquad
\bar z_\ell=\nabla_{z_\ell}L.
\]
For an elementwise activation,
\[
\boxed{
\bar z_\ell=\bar h_\ell\odot \phi_\ell'(z_\ell).
}
\]

\paragraph{Core derivation: affine-layer gradients.}
The affine preactivation is
\[
z=Wh+b.
\]
Its differential is
\[
dz=dW\,h+W\,dh+db.
\]
For upstream adjoint \(\bar z=\nabla_zL\),
\[
dL=\bar z^\top dz
=\bar z^\top dW\,h+\bar z^\top W\,dh+\bar z^\top db.
\]
Using the Frobenius pairing,
\[
\bar z^\top dW\,h
=\operatorname{tr}(\bar z^\top dW h)
=\operatorname{tr}(h\bar z^\top dW)
=\langle \bar z h^\top,dW\rangle_F.
\]
Thus
\[
\boxed{
\nabla_W L=\bar z h^\top,
\qquad
\nabla_b L=\bar z,
\qquad
\bar h=W^\top\bar z.
}
\]
Layerwise backpropagation repeats these formulas from \(\ell=L\) down to \(1\).

\paragraph{Algorithmic form.}
Forward pass:
\[
h_0=x,\qquad
z_\ell=W_\ell h_{\ell-1}+b_\ell,\qquad
h_\ell=\phi_\ell(z_\ell).
\]
Initialize
\[
\bar h_L=\nabla_{h_L}L.
\]
For \(\ell=L,L-1,\ldots,1\):
\[
\bar z_\ell=\bar h_\ell\odot\phi_\ell'(z_\ell),
\]
\[
\nabla_{W_\ell}L=\bar z_\ell h_{\ell-1}^\top,
\qquad
\nabla_{b_\ell}L=\bar z_\ell,
\]
\[
\bar h_{\ell-1}=W_\ell^\top\bar z_\ell.
\]
For minibatches, gradients are summed or averaged over examples. In row-batch notation, if \(H_{\ell-1}\in\mathbb R^{B\times d_{\ell-1}}\) and \(\bar Z_\ell\in\mathbb R^{B\times d_\ell}\), then
\[
\nabla_{W_\ell}L=\bar Z_\ell^\top H_{\ell-1}
\]
up to the chosen sum or average convention.

\paragraph{Concrete example.}
Let a one-layer linear network be
\[
\widehat y=wx+b,\qquad L=\frac12(\widehat y-y)^2.
\]
Here \(z=\widehat y\), \(\bar z=\widehat y-y\). Therefore
\[
\frac{\partial L}{\partial w}=(\widehat y-y)x,\qquad
\frac{\partial L}{\partial b}=\widehat y-y.
\]
For \(x=2\), \(y=5\), \(w=1\), \(b=1\), the prediction is \(3\), so \(\bar z=-2\). The gradients are
\[
\partial L/\partial w=-4,\qquad
\partial L/\partial b=-2.
\]
A gradient step with positive learning rate increases both \(w\) and \(b\), raising the prediction toward \(5\).

\paragraph{Computational neuroscience example.}
A differentiable model may map stimulus history through a filter and nonlinearity to a predicted firing rate. Backpropagation computes how the spike-prediction loss changes with each filter coefficient. In a recurrent neural model, backpropagation through time unrolls the recurrence and applies the same reverse-mode logic across time steps.

\paragraph{AI and machine-learning example.}
Training a multilayer perceptron, CNN, transformer, or value network uses this same pattern. Libraries implement local vector-Jacobian products for matrix multiplication, convolution, normalization, softmax, and attention. The user sees \(\nabla_\theta L\), but internally each module contributes local backward rules.

\paragraph{Common misconceptions and failure modes.}
Backpropagation is not finite differences; it does not estimate gradients by perturbing each parameter separately. The derivative of the activation uses the cached preactivation \(z_\ell\), not the next layer's value. Gradients for weights are outer products of an output-side adjoint and an input-side activation. Forgetting to average or sum consistently over a minibatch changes the effective learning rate.

\paragraph{Exercises.}
\begin{enumerate}[label=\textbf{38.2.\arabic*.}, leftmargin=*]
\item \textbf{Mechanical.} If \(\bar z_\ell\in\mathbb R^{10}\) and \(h_{\ell-1}\in\mathbb R^{4}\), state the shape of \(\nabla_{W_\ell}L=\bar z_\ell h_{\ell-1}^\top\).
\item \textbf{Proof or derivation.} Use \(dL=\bar z^\top(dW h+Wdh+db)\) to derive \(\nabla_WL\), \(\nabla_bL\), and \(\bar h\).
\item \textbf{Numerical/computational.} For \(\widehat y=wx+b\), \(L=\frac12(\widehat y-y)^2\), compute gradients at \(x=3\), \(y=1\), \(w=2\), \(b=-1\).
\item \textbf{Interpretation.} Why must the forward pass cache \(h_{\ell-1}\) and \(z_\ell\) for each layer?
\item \textbf{Misconception/debugging.} A student computes \(\nabla_WL=h\bar z^\top\) for \(W\in\mathbb R^{d_{\mathrm{out}}\times d_{\mathrm{in}}}\). Use shapes to show why the correct column-vector formula is \(\bar z h^\top\).
\end{enumerate}

\subsection{38.3 Differentiable programming}

\paragraph{Guiding question.}
What changes when the object being differentiated is not a simple layer stack but an arbitrary numerical program?

\paragraph{Beginner-facing intuition.}
Differentiable programming treats a program as a parameterized function that can be optimized by gradients. The program might contain neural networks, physics simulators, ODE solvers, rendering steps, control loops, probabilistic losses, or iterative algorithms. If the operations are differentiable, automatic differentiation can propagate gradients through the executed computation.

This generalizes backpropagation. A neural network is one differentiable program. A differentiable simulator plus a neural policy plus a loss is another.

\paragraph{Formal definitions.}
Let a program \(P_\theta\) map input \(x\) to output \(u\):
\[
u=P_\theta(x).
\]
Let a scalar objective be
\[
L(\theta)=\ell(P_\theta(x),y).
\]
During execution, the program creates a computational graph of primitive operations. Reverse-mode automatic differentiation computes
\[
\nabla_\theta L
\]
by applying vector-Jacobian products through those operations in reverse order.

Programs may include control flow. A fixed execution trace with a finite number of operations is differentiable if the executed primitive operations are differentiable at the encountered values. Branching decisions and discrete indexing can create nondifferentiabilities or zero gradients with respect to the values that control them.

\paragraph{Core mechanism: differentiating through an iterative update.}
Consider \(K\) steps of gradient descent inside a program:
\[
u_{k+1}=u_k-\eta\nabla_u E_\theta(u_k),
\qquad k=0,\ldots,K-1,
\]
and final loss
\[
L=\ell(u_K).
\]
For a fixed \(K\), this is a composition of differentiable maps if \(E_\theta\) is smooth. Backpropagation through the iterations applies
\[
\bar u_k=
\left(\frac{\partial u_{k+1}}{\partial u_k}\right)^\top\bar u_{k+1}.
\]
Since
\[
\frac{\partial u_{k+1}}{\partial u_k}
=I-\eta\nabla_u^2E_\theta(u_k),
\]
the backward pass depends on Hessian-vector products. Differentiable programming systems can compute these products without forming full Hessian matrices.

\paragraph{Concrete example.}
Let
\[
u_1=\theta x,\qquad u_2=u_1^2,\qquad L=\frac12(u_2-y)^2.
\]
At \(x=2\), \(\theta=3\), and \(y=40\),
\[
u_1=6,\qquad u_2=36,\qquad L=8.
\]
The backward values are
\[
\bar u_2=u_2-y=-4,
\qquad
\bar u_1=\bar u_2(2u_1)=-48,
\]
and
\[
\frac{\partial L}{\partial \theta}=\bar u_1 x=-96.
\]
This is the same graph logic whether the nodes are called layers, simulation steps, or program variables.

\paragraph{Computational neuroscience example.}
A differentiable model of behavior might simulate a latent decision variable, pass it through a softmax choice rule, and compare predicted choices with data. Gradients can fit drift rates, noise scales, or policy parameters. If the simulator contains a hard threshold for spike generation or decision commitment, a differentiable surrogate or likelihood-based treatment may be needed.

\paragraph{AI and machine-learning example.}
Differentiable rendering fits scene parameters from images. Differentiable physics tunes controller parameters through a simulator. Neural ODEs differentiate through continuous-time dynamics. Meta-learning differentiates through an inner training loop. In each case, the model is a program whose intermediate values form the graph.

\paragraph{Common misconceptions and failure modes.}
Automatic differentiation differentiates the executed computation, not an idealized mathematical intention that was not in the code. Discrete choices, rounding, argmax, sampling, and data-dependent indexing can break useful gradients. Differentiability does not imply numerical stability. Memory can be a bottleneck because reverse mode needs forward values; checkpointing trades recomputation for memory. Custom gradients must match the intended mathematical derivative or be clearly treated as surrogate gradients.

\paragraph{Exercises.}
\begin{enumerate}[label=\textbf{38.3.\arabic*.}, leftmargin=*]
\item \textbf{Mechanical.} Name three primitive operations in a differentiable program and one operation that is typically non-differentiable.
\item \textbf{Proof or derivation.} For \(u_{k+1}=u_k-\eta\nabla E(u_k)\), derive \(\partial u_{k+1}/\partial u_k=I-\eta\nabla^2E(u_k)\).
\item \textbf{Numerical/computational.} For \(u_1=\theta x\), \(u_2=u_1^2\), \(L=\frac12(u_2-y)^2\), compute \(\partial L/\partial\theta\) at \(x=1\), \(\theta=2\), \(y=5\).
\item \textbf{Interpretation.} Why might a hard spike threshold cause trouble for gradient-based fitting?
\item \textbf{Misconception/debugging.} A student says, ``If a simulator runs in an autodiff library, the gradients must be meaningful.'' Explain why branches, discontinuities, and numerical error still need inspection.
\end{enumerate}

\subsection{38.4 Gradient pathologies}

\paragraph{Guiding question.}
Why can mathematically correct gradients become too small, too large, or misleading for learning?

\paragraph{Beginner-facing intuition.}
Backpropagation multiplies local derivative factors. If many factors have norms less than \(1\), gradients can vanish. If many have norms greater than \(1\), gradients can explode. Activations can saturate, ReLUs can die, and poorly conditioned objectives can make gradient descent zigzag or crawl. These are not mysterious failures; they are consequences of the chain rule and numerical optimization.

\paragraph{Formal definitions.}
For a composition
\[
h_L=F_L(F_{L-1}(\cdots F_1(x))),
\]
the input gradient of a scalar loss is
\[
\nabla_xL
=J_{F_1}(x)^\top J_{F_2}(h_1)^\top\cdots J_{F_L}(h_{L-1})^\top\nabla_{h_L}L.
\]
Taking norms,
\[
\|\nabla_xL\|
\leq
\left(\prod_{\ell=1}^L \|J_{F_\ell}\|\right)\|\nabla_{h_L}L\|.
\]
This bound shows how repeated Jacobian factors can shrink or grow gradients exponentially with depth or time.

\paragraph{Core mechanisms.}
\emph{Saturation.} The sigmoid
\[
\sigma(u)=\frac1{1+e^{-u}}
\]
has derivative
\[
\sigma'(u)=\sigma(u)(1-\sigma(u)).
\]
When \(u\) is very positive or very negative, \(\sigma(u)\) is near \(1\) or \(0\), so \(\sigma'(u)\) is near \(0\). Gradients through saturated units are small.

\emph{Exploding and vanishing in repeated linear maps.} If
\[
h_t=Ah_{t-1},
\]
then
\[
h_T=A^T h_0
\]
and backward gradients contain \((A^\top)^T\). Eigenvalues or singular values below \(1\) shrink signals; those above \(1\) amplify them.

\emph{Dead ReLUs.} ReLU derivative is \(0\) for negative preactivation and \(1\) for positive preactivation, ignoring the kink at \(0\). If a unit's preactivation is negative for all examples, its incoming weights may receive no gradient from that unit.

\paragraph{Concrete example.}
Suppose a scalar recurrent derivative is \(0.5\) at each of \(20\) time steps. The gradient factor is
\[
0.5^{20}\approx9.5\times10^{-7}.
\]
If the derivative is \(1.2\) at each step, the factor is
\[
1.2^{20}\approx38.3.
\]
The first gradient nearly vanishes; the second can become large enough to destabilize updates.

\paragraph{Computational neuroscience example.}
Fitting recurrent neural models to long behavioral or neural sequences requires gradients through time. If recurrent dynamics are strongly contracting, early inputs or parameters may receive tiny gradients. If dynamics are unstable, gradients can explode. Biological recurrent circuits also face stability and sensitivity tradeoffs, but that does not mean they solve them with backpropagation through time.

\paragraph{AI and machine-learning example.}
Residual connections provide shorter gradient paths:
\[
h_{\ell+1}=h_\ell+F_\ell(h_\ell).
\]
The derivative contains
\[
I+J_{F_\ell},
\]
so there is a direct identity component. Normalization, careful initialization, gating, gradient clipping, and adaptive optimizers are also used to manage gradient scale and conditioning.

\paragraph{Common misconceptions and failure modes.}
Small training progress is not always caused by vanishing gradients; data, loss scaling, bugs, or learning-rate choice can be responsible. Exploding gradients are not fixed by merely lowering the learning rate if the model dynamics remain unstable. ReLU units with zero output on one batch are not necessarily dead forever, but persistently inactive units are suspicious. Gradient clipping controls update magnitude but does not solve all modeling or credit-assignment problems.

\paragraph{Gradient debugging checklist.}
Useful checks include monitoring gradient norms by layer, checking activation statistics, verifying loss scale, comparing autodiff gradients with finite differences on a tiny model, testing one-batch overfitting, confirming train/eval mode for normalization layers, and inspecting whether parameters that should receive gradients have nonzero gradients.

\paragraph{Exercises.}
\begin{enumerate}[label=\textbf{38.4.\arabic*.}, leftmargin=*]
\item \textbf{Mechanical.} Compute \(0.8^{10}\) and \(1.1^{10}\). Which corresponds to shrinking gradient factors and which to growing factors?
\item \textbf{Proof or derivation.} Use the submultiplicative norm inequality to derive \(\|\nabla_xL\|\leq(\prod_\ell\|J_{F_\ell}\|)\|\nabla_{h_L}L\|\).
\item \textbf{Numerical/computational.} Evaluate \(\sigma'(u)=\sigma(u)(1-\sigma(u))\) approximately at \(u=0\) and \(u=6\). Explain the difference.
\item \textbf{Interpretation.} Why can residual connections help gradient flow in very deep networks?
\item \textbf{Misconception/debugging.} A model's first layer has exactly zero gradients for many batches. List two possible causes and one diagnostic check.
\end{enumerate}

Backpropagation is the operational form of the chain rule for modern learning systems. It turns the parameterized functions of Chapter 37 into trainable models, and differentiable programming extends the same principle to simulations, control loops, and scientific models. Later chapters study how this gradient machinery behaves in deep optimization, specialized architectures, generative models, and computational neuroscience applications.

\clearpage
